\begin{lemma}[Interchangeability for a closed Tashkinov tree]
Let $(G,e_0,\varphi)$ be a $k$-triple, and let $T$ be a Tashkinov
tree with respect to $e_0$ and $\varphi$ whose vertex set is closed under
$\varphi$.  If $\alpha$ or $\beta$ is missing on $T$, then $\alpha$ and
$\beta$ are $T$-interchangeable under $\varphi$; equivalently, at most one
nontrivial $(\alpha,\beta)$-path component meets $V(T)$.
\end{lemma}
\begin{proof}
We use path component in the precise sense of
\noderef{TwoColorPathComponent}, and interchangeability in the sense of
\noderef{TreeInterchangeable}.  The phrase ``missing on $T$'' means the
union \noderef{TreeMissingColors} over the vertex set
\noderef{TreeSequenceVertexSet}.  If $\alpha=\beta$, the former definition
admits no path component, because it requires two distinct colors, and the
conclusion is vacuous.  Hence assume $\alpha\ne\beta$.

Suppose for a contradiction that two different $(\alpha,\beta)$-path
components $Q_1,Q_2$ in the sense of \noderef{TwoColorComponent} meet
$V(T)$.  By
\noderef{TwoColorComponentsDisjoint}, they are disjoint.  The properness
clause of the $k$-triple in \noderef{KTriple}, together with
\noderef{TwoColorIncidenceDegreeAtMostTwo}, says that every node of either
component has incidence-graph degree at most two.  Each component contains
an edge-node and has a node of degree at most one by
\noderef{TwoColorPathComponent}.  The spanning path supplied by
\noderef{TwoColorComponentSpanningPath} uses every node of the component.
Because the component also meets $V(T)$, its stipulated edge-node is not an
isolated node, so this spanning path is nontrivial.
No additional component edge can join nonconsecutive nodes of this path:
such an edge either gives an internal path node a third neighbor, or joins
the two path ends and makes every node have degree two, both contrary to
the preceding degree facts.  Thus the incidence component itself is a
finite nontrivial simple path.
Both ends are vertex-nodes.  Indeed, by \noderef{Incident} and
\noderef{TwoColorIncidenceGraph}, the neighbors of a colored edge-node are
exactly the vertex-nodes represented by the entries of its endpoint pair.
The loopless law in \noderef{LooplessMultigraph} makes those two entries
distinct, so every colored edge-node has degree two and cannot be an end of
the component path.  At each vertex-node end at least one of
$\alpha,\beta$ is missing, by the definitions of
\noderef{PresentColors} and \noderef{MissingColors}.

The hypotheses of \noderef{TashTree} are exactly the $k$-triple and the
Tashkinov-tree hypotheses here.  Consequently $V(T)$ is elementary under
$\varphi$ in the sense of \noderef{ColoredVertexSetElementary}.  Thus at
most one vertex of $T$ misses $\alpha$, and at most one misses $\beta$.
In particular, among the four ends of the two disjoint paths, at most two
can lie in $V(T)$.

In fact exactly one of $\alpha,\beta$ is missing on $T$.  At least one is
missing by hypothesis.  If both were missing, closedness in the sense of
\noderef{ColoredVertexSetClosed} would forbid every $\alpha$- or
$\beta$-colored edge in \noderef{BoundaryEdges}.  A two-color component meeting $T$ could
then contain no vertex outside $T$: along a component path, the first edge
leaving $V(T)$ would be just such a boundary edge.  Hence all four ends of
$Q_1,Q_2$ would lie in $V(T)$ and miss one of the two colors, contrary to
elementariness.  Interchanging the names if necessary, assume from now on
that
\[
 \alpha\in\overline\varphi(T),\qquad
 \beta\notin\overline\varphi(T).
\]
Closedness now implies that every boundary edge of either $Q_i$ has color
$\beta$.  Since no vertex of $T$ misses $\beta$ and at most one vertex of
$T$ misses $\alpha$, at most one of the four path ends lies in $T$.
Therefore at least three ends lie outside $T$.

Starting at three distinct outside ends and traversing their component
paths toward their first vertex of $T$ gives three pairwise edge-disjoint
exit paths $P_1,P_2,P_3$.  Write $a_i$ for the outside end, $b_i$ for the
first tree vertex, and $f_i$ for the last multigraph edge, incident with
$b_i$.  Each $P_i$ is a simple incidence path of length at least two, its
only tree vertex is $b_i$, and one of $\alpha,\beta$ is missing at $a_i$;
these are exactly the conditions of \noderef{TreeExitPath}.  Moreover
$\varphi(f_i)=\beta$ by closedness.  The $b_i$ are distinct.  Indeed, if
$b_i=b_j$ for $i\ne j$, edge-disjointness gives $f_i\ne f_j$, while
\noderef{Incident} says both edges are incident with that common vertex.
They are therefore adjacent by \noderef{AdjacentEdges}.  Since both have
color $\beta$, this contradicts properness in the sense of
\noderef{IsProperEdgeColoring}.
The order encoded by \noderef{TreeVertexOrder} is linear on the distinct
displayed vertices, so relabel the paths so that
$b_1\prec_T b_2\prec_T b_3$.

Call any tuple with all the properties just listed a counterexample.  More
explicitly, its coloring is part of a $k$-triple on the fixed
$(G,k,e_0)$, its tree is closed and Tashkinov for that coloring, at least
one of its two colors is missing on the tree, and its three exit paths have distinct
outside and tree ends, are the initial-to-first-tree portions of their full
two-color component paths, and are pairwise edge-disjoint.  The preceding
paragraph constructs such a tuple from the assumed failure.  The set of
the three-path total lengths is therefore a nonempty set of natural
numbers; choose a counterexample minimizing
$|P_1|+|P_2|+|P_3|$.  Here and below $|P_i|$ is half the incidence-graph
length, namely the number of multigraph edge-nodes of $P_i$, so strict
shortening is the same in either convention.

Every vertex misses a color.  Indeed, the $k$-triple gives
$k\geq\Delta(G)+1$.  For a vertex $v$, choose for each color in
\noderef{PresentColors} an incident edge carrying that color.  This choice
is injective, since one edge has only one coloring value.  Hence the number
of present colors is at most the number of incident edge identities, which
is $d_G(v)$ by \noderef{Degree}; and \noderef{MaximumDegree} gives
$d_G(v)\leq\Delta(G)<k$.  The complement definition
\noderef{MissingColors} therefore leaves a missing color.  Choose
$\gamma\in\overline\varphi(b_3)$ and put
\[
 \sigma_1=\varphi/(G-T,\alpha,\gamma)
\]
using \noderef{OutsideTreeColorSwap}.  Both $\alpha$ and $\gamma$ are
missing on $T$, so closedness forbids either color on a boundary edge.
We record all coloring checks for this operation.  Two adjacent edges both
incident with $T$ in the sense of \noderef{TreeIncidentEdge} are fixed.
Two adjacent edges neither incident with
$T$ undergo the same transposition of $\alpha,\gamma$.  If exactly one is
incident with $T$, their common vertex is outside $T$ and the fixed edge
is a boundary edge, which has neither color; swapping the other edge cannot
create a conflict.  Thus $\sigma_1$ is proper.  The distinguished edge
$e_0$ is the first edge of the tree-sequence in \noderef{TreeSequence}, so
it is incident with $T$ and remains uncolored; every other edge remains
colored because the operation only transposes colored values.  Hence the
unique-uncolored clause \noderef{IsOneEdgeUncolored} is preserved.

Every edge incident with $T$ is fixed, so the present and missing color
sets at every tree vertex are unchanged.  Every displayed tree edge and
every earlier missing-color witness in \noderef{TashkinovTree} is therefore
unchanged.  The same observation preserves closedness.  The graph,
chromatic-index, criticality, and degree clauses do not depend on the
coloring.  Consequently $(G,e_0,\sigma_1)$ is again a $k$-triple and $T$
is again a closed Tashkinov tree.

The boundary edges $f_1,f_2,f_3$ keep color $\beta$.  On the portion of
each $P_i$ outside $T$, $\alpha$ is changed to $\gamma$ and $\beta$ is
fixed.  More explicitly, every path edge other than $f_i$ is not incident
with $T$: in the full original component path its two endpoint vertices
occur before the first tree vertex $b_i$.  At the outside end the
corresponding missing status is changed in the same way.  Hence the same
three edge sequences are
$(T,\sigma_1,\{\gamma,\beta\})$-exit paths.  In particular $P_3$ is the
entire $(\gamma,\beta)$-component between $a_3$ and $b_3$: its outside end
still misses one of the two colors, while at $b_3$ the color $\gamma$ is
missing and the unique incident $\beta$-edge of this component is $f_3$.
Because $f_3$ has color $\beta$ at $b_3$, the chosen missing color
$\gamma$ is distinct from $\beta$.

Let $\sigma_2$ be obtained by interchanging $\gamma$ and $\beta$ on the
component $P_3$, that is, by the operation
\noderef{TwoColorComponentSwap} rooted at a node of $P_3$.
Properness follows from \noderef{TwoColorComponentSwapProper}.  This swap
again only transposes colored values, so it preserves the assertion that
$e_0$ is the unique uncolored edge.  It interchanges the two missing
statuses at the two ends of $P_3$ and changes no other incidence; hence
\[
 \beta\in\overline{\sigma_2}(b_3).
\]
The components containing $P_1$ and $P_2$ are disjoint from the component
$P_3$: otherwise that component, whose support is exactly $P_3$, would
contain the distinct tree end $b_1$ or $b_2$.  This also follows from
\noderef{TwoColorComponentsDisjoint}.  Hence their edges and endpoint
missing statuses are unchanged.  Thus all three $P_i$ remain
$(\gamma,\beta)$-exit paths under $\sigma_2$, with the colors merely
interchanged along $P_3$.

Let $q$ be the displayed index of $b_3$.  Since two distinct tree vertices
precede it, $1\le q\le\operatorname{length}(T)$, and
\noderef{TreeSequencePrefix} gives the tree-sequence
$S=T|b_3$.  The swap on $P_3$ meets $T$ only at $b_3$.  Its sole edge
incident with $b_3$ is the boundary edge $f_3$, not a displayed tree edge;
no edge or missing witness at a vertex strictly preceding $b_3$ changes.
Although the missing sets at $b_3$ exchange $\gamma$ and $\beta$, that last
vertex is not an earlier witness for any edge of $S$.  It follows directly
from \noderef{TashkinovTree} that $S$ is a Tashkinov tree under $\sigma_2$.

For clarity, let $x_i$ be the outside endpoint of the boundary edge $f_i$;
this need not be the remote path end $a_i$ when $P_i$ has more than one
edge.  Append first $(f_1,x_1)$ and then $(f_2,x_2)$ to $S$, obtaining
$T'$.  We check that these really are the two augmentations represented by
\noderef{TashkinovAugmentation}.  The vertices $x_1,x_2$ are outside the
original $T$, because $b_i$ was the first tree vertex of $P_i$.  The edge
identities $f_1,f_2$ are not displayed edges of $S$, since each is a
boundary edge while every displayed tree edge has both ends in the tree.
They are distinct because the exit paths were chosen edge-disjoint.  Also
$x_1\ne x_2$: otherwise the distinct edges $f_1,f_2$ would both be
incident with $x_1$ by \noderef{Incident}, hence adjacent by
\noderef{AdjacentEdges}; as both still have color $\beta$ under $\sigma_2$,
this would violate \noderef{IsProperEdgeColoring}.  Each $f_i$
attaches its new outside endpoint $x_i$ to the earlier displayed vertex
$b_i$.  Finally $\sigma_2(f_1)=\sigma_2(f_2)=\beta$, and $\beta$ is
missing at the earlier vertex $b_3$ both before and after the first
augmentation, since neither $f_1$ nor $f_2$ is incident with the distinct
vertex $b_3$.  Thus the old entries agree exactly as required by
\noderef{TreeSequenceInitialSegment}, and the two steps form a
\noderef{TashkinovAugmentationSequence} from $S$ to $T'$.
By \noderef{TashkinovAugmentationSequencePreservesTashkinovTree}, $T'$ is a
Tashkinov tree under $\sigma_2$.

Choose a closure $T''$ of $T'$ under $\sigma_2$ by
\noderef{TashkinovClosureExists}.  By the definition
\noderef{TashkinovClosure}, there is a Tashkinov augmentation sequence from
$T'$ to $T''$ and $V(T'')$ is closed.  A second application of
\noderef{TashkinovAugmentationSequencePreservesTashkinovTree} shows that
$T''$ is a Tashkinov tree.  At every augmentation, the old vertices and
edges remain entrywise fixed by \noderef{TreeSequenceInitialSegment}; hence
$V(S)\subseteq V(T')\subseteq V(T'')$.  All coloring-independent clauses of the
$k$-triple were already preserved, so \noderef{TashTree} now gives that
$V(T'')$ is elementary under $\sigma_2$.

We claim that none of $a_1,a_2,a_3$ lies in $V(T'')$.  Suppose
$a_i\in V(T'')$.  At least one of $\beta,\gamma$ is missing at $a_i$ under
$\sigma_2$, while $\beta$ is missing at the distinct vertex
$b_3\in V(S)\subseteq V(T'')$.  Elementariness therefore forces
$\gamma$ to be missing at $a_i$.  Both $\beta$ and $\gamma$ are now
missing somewhere on $T''$.  Closedness of $T''$ says that no boundary
edge has either color.

Each $b_j$ belongs to the prefix $S$ and hence to $T''$.  Traverse $P_j$
from $b_j$ toward $a_j$.  Its edges alternate $\beta,\gamma$ under
$\sigma_2$.  If the next vertex ever lay outside $V(T'')$, the first edge
crossing from inside to outside would be a $\beta$- or $\gamma$-colored
boundary edge, contradicting closedness.  Induction along the finite path
therefore puts all three paths, and in particular all three distinct
$a_j$, inside $V(T'')$.  At every $a_j$ at least one of
$\beta,\gamma$ is missing.  Since $\beta$ is already missing at the
distinct vertex $b_3$, elementariness forces $\gamma$ to be missing at
each $a_j$.  This gives the same missing color at three distinct vertices
of $V(T'')$, contradicting elementariness.  The claim follows.

For each $i$, traverse $P_i$ from the outside vertex $a_i$ and stop at its
first vertex $c_i$ in $V(T'')$; such a vertex exists because
$b_i\in V(S)\subseteq V(T'')$.  Call the resulting initial subpath $L_i$,
and call its final multigraph edge $g_i$.  The preceding claim keeps
$a_i$ outside $V(T'')$, while the definition of $c_i$ makes it the only
vertex-node of $L_i$ in $V(T'')$.  Thus $L_i$ is nontrivial.  It remains a
simple path in the $(\gamma,\beta)$-incidence graph, and one of the two
colors is missing at $a_i$.  These facts verify every clause of
\noderef{TreeExitPath}, so $L_i$ is a $T''$-exit path.  The $a_i$ remain
distinct, and the $L_i$ remain pairwise edge-disjoint because they are
subpaths of the corresponding $P_i$.

We next determine the colors of the three final edges.  The two consecutive
vertex-nodes around $g_i$ are the vertex just before $c_i$, which lies
outside $V(T'')$, and $c_i$, which lies inside.  Adjacency in
\noderef{TwoColorIncidenceGraph} and the incidence definition
\noderef{Incident} show that these are the two ends of $g_i$; they are
distinct by the loopless law in \noderef{LooplessMultigraph}.  Hence
$g_i\in\partial_G(V(T''))$ by \noderef{BoundaryEdges}.  The same
two-color adjacency says that $\sigma_2(g_i)$ is either $\gamma$ or
$\beta$.  Now $\beta$ is missing at $b_3\in V(T'')$, so it is missing on
$T''$ by \noderef{TreeMissingColors}.  Since $V(T'')$ is closed,
\noderef{ColoredVertexSetClosed} forbids color $\beta$ on every boundary
edge.  Consequently
\[
 \sigma_2(g_i)=\gamma\qquad(i=1,2,3).
\]

The new tree ends $c_1,c_2,c_3$ are distinct.  If $c_i=c_j$ for
$i\ne j$, pairwise edge-disjointness gives $g_i\ne g_j$ and
\noderef{Incident} makes both edges incident with that common vertex.
Thus \noderef{AdjacentEdges} makes them adjacent.  Their common color
$\gamma$ then contradicts properness as defined in
\noderef{IsProperEdgeColoring}.

It remains to verify the full-component-prefix condition, rather than only
the exit-path condition.  Fix $i$ and let $C_i$ be the full
$(\gamma,\beta)$-component through the vertex-node $a_i$, as defined by
\noderef{TwoColorComponent}.  Missingness at $a_i$ means, by
\noderef{MissingColors} and \noderef{PresentColors}, that no incident edge
of the missing color is adjacent to $a_i$ in
\noderef{TwoColorIncidenceGraph}.  There is at most one incident edge of
the other color: two distinct such identities would both be incident with
$a_i$ by \noderef{Incident}, hence adjacent by \noderef{AdjacentEdges},
and their equal colors would contradict
\noderef{IsProperEdgeColoring}.  Thus the incidence-graph degree of $a_i$
is at most one.  The first edge-node of the nontrivial path $L_i$ is a
neighbor of $a_i$, so the degree is exactly one.

As recorded above, $\gamma\ne\beta$: under $\sigma_1$ the color $\gamma$
is missing at $b_3$, while its incident edge $f_3$ has color $\beta$.
The component $C_i$ contains the edge-node $g_i$ and has the degree-one
node $a_i$, so it is a path component in the sense of
\noderef{TwoColorPathComponent}.  Apply
\noderef{TwoColorComponentSpanningPath} to the proper coloring $\sigma_2$
and orient the resulting simple spanning path $R_i$ from $a_i$.  This
orientation is possible because $a_i$ cannot be internal in $R_i$: an
internal occurrence would give it two distinct neighbors, whereas it has
degree one.

We compare $R_i$ with $L_i$ one node at a time.  Their first node is
$a_i$, and its unique neighbor forces their next nodes to agree.  Suppose
inductively that they agree through some internal node $z$ of $L_i$.
Because $L_i$ is simple, its predecessor and successor at $z$ are two
distinct neighbors.  The bound
\noderef{TwoColorIncidenceDegreeAtMostTwo} says that $z$ has no third
neighbor.  The spanning path $R_i$ has arrived through the common
predecessor and cannot repeat it, so its next node must be the displayed
successor of $L_i$.  Induction continues up to $c_i$.  Therefore $L_i$ is
exactly the initial segment of the full component path $R_i$ from its end
$a_i$ through $c_i$.  Since by construction no earlier vertex-node of
$L_i$ belongs to $V(T'')$, it is precisely the initial-to-first-$T''$
prefix required in the counterexample class.  This argument applies to
all three indices.

Both ends $b_1,x_1$ of $f_1$ already belong to $T'$, hence to $T''$.
Consequently the first-entry subpath $L_1$ stops no later than $x_1$ and
does not contain $f_1$, whereas $P_1$ does.  Thus
\[
 |L_1|+|L_2|+|L_3|<|P_1|+|P_2|+|P_3|.
\]
Finally relabel the three $L_i$ if necessary so that their now distinct
tree ends satisfy $c_1\prec_{T''}c_2\prec_{T''}c_3$ in the linear order
\noderef{TreeVertexOrder}; this permutation leaves their total length
unchanged.  We can now restate every counterexample condition for
$(\sigma_2,T'',\gamma,\beta,L_1,L_2,L_3)$.  The graph, $k$, and
distinguished edge are the fixed ones; $(G,e_0,\sigma_2)$ is a $k$-triple;
$T''$ is a closed Tashkinov tree; $\gamma\ne\beta$ and $\beta$ is missing
on $T''$; each $L_i$ is a nontrivial simple $(\gamma,\beta)$-exit path
from the distinct outside end $a_i$ to the distinct ordered tree end
$c_i$; each is exactly the initial-to-first-tree prefix of its full
two-color component path; and the three paths are pairwise edge-disjoint.
Thus this tuple belongs to the very counterexample class minimized above.
Its strictly smaller total length contradicts the minimum choice.
Therefore two different path components cannot both meet $T$, which is
exactly the asserted $T$-interchangeability.
\end{proof}
